\documentclass[12pt,a4paper]{article}
\usepackage[UTF8]{ctex}
\usepackage{geometry}
\usepackage{listings}
\usepackage{xcolor}
\usepackage{hyperref}
\usepackage{booktabs}
\usepackage{graphicx}
\usepackage{amsmath}
\usepackage{float}
\usepackage{tcolorbox}
\usepackage{enumitem}
\usepackage{amssymb}

% 页面设置
\geometry{left=2.5cm,right=2.5cm,top=2.5cm,bottom=2.5cm}

% 代码高亮设置
\lstset{
    language=Python,
    basicstyle=\ttfamily\small,
    keywordstyle=\color{blue}\bfseries,
    commentstyle=\color{gray},
    stringstyle=\color{orange},
    numbers=left,
    numberstyle=\tiny\color{gray},
    stepnumber=1,
    numbersep=8pt,
    backgroundcolor=\color{gray!10},
    frame=single,
    frameround=fttt,
    breaklines=true,
    showstringspaces=false,
    tabsize=4,
    captionpos=b,
    xleftmargin=2em,
    xrightmargin=1em,
    aboveskip=1em,
    belowskip=1em
}

% 知识点框
\tcbuselibrary{skins,breakable}
\newtcolorbox{knowledgebox}[1][]{
    colback=blue!5!white,
    colframe=blue!50!black,
    fonttitle=\bfseries,
    title=#1,
    breakable
}

% 练习框
\newtcolorbox{practicebox}[1][]{
    colback=green!5!white,
    colframe=green!50!black,
    fonttitle=\bfseries,
    title=#1,
    breakable
}

% 注意框
\newtcolorbox{tipbox}[1][]{
    colback=yellow!10!white,
    colframe=orange!70!black,
    fonttitle=\bfseries,
    title=#1,
    breakable
}

% 超链接设置
\hypersetup{
    colorlinks=true,
    linkcolor=blue,
    filecolor=magenta,      
    urlcolor=cyan,
    bookmarks=true,
    pdfstartview=FitH
}

% 标题设置
\title{\textbf{Pandas入门与数据处理实战手册}\\[0.5em]\large 零基础入门版}
\author{竞赛培训组}
\date{\today}

\begin{document}

\maketitle
\tableofcontents
\newpage

%========================================
\section{Python基础入门}
%========================================

\subsection{什么是Python？}

\begin{knowledgebox}[知识点：Python是什么]
Python是一种\textbf{编程语言}，就像中文、英语是人类交流的语言一样，Python是人和计算机交流的语言。

\textbf{Python的特点：}
\begin{itemize}
    \item \textbf{简单易学}：语法接近自然语言，适合初学者
    \item \textbf{功能强大}：可以做数据分析、网站开发、人工智能等
    \item \textbf{应用广泛}：大数据、机器学习、自动化办公等领域都在用
\end{itemize}

\textbf{为什么要学Python？}
\begin{itemize}
    \item 数据处理必备技能
    \item 竞赛中必须使用的工具
    \item 就业市场需求大
\end{itemize}
\end{knowledgebox}

\subsection{安装Python}

在使用Pandas之前，需要先安装Python。

\textbf{步骤：}
\begin{enumerate}
    \item 访问官网：\url{https://www.python.org/downloads/}
    \item 下载最新版本的Python（建议3.8或以上）
    \item 运行安装程序，\textbf{一定要勾选"Add Python to PATH"}
    \item 点击"Install Now"完成安装
\end{enumerate}

\begin{tipbox}[安装注意事项]
\begin{itemize}
    \item 安装时务必勾选"Add Python to PATH"，否则后面会有问题
    \item 如果已经安装了Python，可以跳过这一步
    \item 安装完成后，重启电脑确保环境变量生效
\end{itemize}
\end{tipbox}

\subsection{验证Python安装}

安装完成后，验证一下是否安装成功：

\begin{enumerate}
    \item 按\texttt{Win + R}键，输入\texttt{cmd}，回车打开命令行
    \item 输入以下命令：
\begin{verbatim}
python --version
\end{verbatim}
    \item 如果显示版本号（如\texttt{Python 3.9.0}），说明安装成功
\end{enumerate}

\subsection{Python基础语法}

\subsubsection{变量和数据类型}

\begin{knowledgebox}[知识点：什么是变量？]
变量就像是给数据起个名字，方便我们以后使用。

\textbf{类比理解：}
\begin{itemize}
    \item 变量名 = 标签
    \item 变量值 = 盒子里的东西
\end{itemize}

比如：\texttt{age = 20}，就是把数字20放进一个叫\texttt{age}的盒子里。
\end{knowledgebox}

\begin{lstlisting}[caption={变量示例}]
# 定义变量（给数据起名字）
name = "张三"      # 字符串（文字）
age = 20          # 整数（数字）
height = 1.75     # 浮点数（小数）
is_student = True # 布尔值（真/假）

# 打印变量
print(name)
print(age)
print(height)
print(is_student)
\end{lstlisting}

\textbf{Python常用数据类型：}

\begin{table}[H]
\centering
\caption{Python数据类型}
\begin{tabular}{|l|l|l|}
\hline
\textbf{类型} & \textbf{说明} & \textbf{示例} \\
\hline
int & 整数 & 20, 100, -5 \\
\hline
float & 浮点数（小数） & 1.75, 3.14 \\
\hline
str & 字符串（文字） & "张三", "hello" \\
\hline
bool & 布尔值（真/假） & True, False \\
\hline
list & 列表（一组数据） & [1, 2, 3] \\
\hline
dict & 字典（键值对） & \{"name": "张三"\} \\
\hline
\end{tabular}
\end{table}

\subsubsection{列表（List）}

\begin{knowledgebox}[知识点：什么是列表？]
列表就是一组数据的集合，用方括号\texttt{[]}表示。

\textbf{类比理解：}
\begin{itemize}
    \item 列表 = 购物清单、成绩单、通讯录
    \item 可以存放多个数据
    \item 数据之间有顺序
\end{itemize}
\end{knowledgebox}

\begin{lstlisting}[caption={列表示例}]
# 创建一个列表
scores = [85, 92, 78, 90, 88]
names = ["张三", "李四", "王五"]

# 访问列表元素（索引从0开始）
print(scores[0])    # 第一个元素：85
print(scores[1])    # 第二个元素：92
print(names[2])     # 第三个元素：王五

# 获取列表长度
print(len(scores))  # 5

# 添加元素
scores.append(95)   # 在末尾添加95
\end{lstlisting}

\subsubsection{字典（Dictionary）}

\begin{knowledgebox}[知识点：什么是字典？]
字典是一种键值对的数据结构，用花括号\texttt{\{\}}表示。

\textbf{类比理解：}
\begin{itemize}
    \item 字典 = 真实字典（查词找释义）
    \item 键（key）= 要查的词
    \item 值（value）= 词的释义
\end{itemize}
\end{knowledgebox}

\begin{lstlisting}[caption={字典示例}]
# 创建一个字典
student = {
    "name": "张三",
    "age": 20,
    "score": 85
}

# 访问字典元素
print(student["name"])   # 张三
print(student["age"])    # 20

# 修改值
student["score"] = 90

# 添加新键值对
student["gender"] = "男"
\end{lstlisting}

\subsubsection{条件判断（if语句）}

\begin{lstlisting}[caption={条件判断示例}]
# 条件判断
age = 18

if age >= 18:
    print("成年人")
else:
    print("未成年人")

# 多条件判断
score = 85

if score >= 90:
    print("优秀")
elif score >= 80:
    print("良好")
elif score >= 60:
    print("及格")
else:
    print("不及格")
\end{lstlisting}

\begin{knowledgebox}[知识点：条件判断的语法]
\begin{itemize}
    \item \texttt{if}：如果条件成立
    \item \texttt{elif}：否则如果（可以有多个）
    \item \texttt{else}：否则（前面都不成立时执行）
    \item \textbf{注意}：Python用\textbf{缩进}表示代码块，不是大括号
\end{itemize}
\end{knowledgebox}

\subsubsection{循环（for循环）}

\begin{lstlisting}[caption={循环示例}]
# 遍历列表
scores = [85, 92, 78, 90]

for score in scores:
    print(score)

# 遍历范围
for i in range(5):      # 0, 1, 2, 3, 4
    print(i)

# 遍历字典
student = {"name": "张三", "age": 20}

for key, value in student.items():
    print(f"{key}: {value}")
\end{lstlisting}

\subsubsection{函数（Function）}

\begin{knowledgebox}[知识点：什么是函数？]
函数就是一段可以重复使用的代码块。

\textbf{类比理解：}
\begin{itemize}
    \item 函数 = 洗衣机、电饭煲（封装好的功能）
    \item 输入 = 脏衣服、米（参数）
    \item 输出 = 干净衣服、米饭（返回值）
\end{itemize}

\textbf{好处：}
\begin{itemize}
    \item 避免重复写代码
    \item 代码更清晰
    \item 便于维护
\end{itemize}
\end{knowledgebox}

\begin{lstlisting}[caption={函数示例}]
# 定义函数
def greet(name):
    """打招呼函数"""
    return f"你好，{name}！"

# 调用函数
result = greet("张三")
print(result)    # 你好，张三！

# 带多个参数的函数
def calculate_total(price, quantity):
    """计算总价"""
    return price * quantity

# 调用函数
total = calculate_total(100, 3)
print(total)     # 300
\end{lstlisting}

\begin{practicebox}[动手练习：Python基础]
完成以下练习，巩固Python基础：

\textbf{练习1：变量和类型}
\begin{lstlisting}
# 创建变量
name = "你的名字"
age = 你的年龄
height = 你的身高

# 打印这些信息
print(f"姓名：{name}")
print(f"年龄：{age}")
print(f"身高：{height}米")
\end{lstlisting}

\textbf{练习2：列表操作}
\begin{lstlisting}
# 创建一个成绩列表
scores = [78, 85, 92, 67, 88]

# 计算平均分
total = sum(scores)
count = len(scores)
average = total / count

print(f"总分：{total}")
print(f"平均分：{average}")
\end{lstlisting}

\textbf{练习3：条件判断}
\begin{lstlisting}
# 判断成绩等级
score = 85

if score >= 90:
    level = "优秀"
elif score >= 80:
    level = "良好"
elif score >= 60:
    level = "及格"
else:
    level = "不及格"

print(f"成绩等级：{level}")
\end{lstlisting}
\end{practicebox}

%========================================
\section{准备工作：安装与导入}
%========================================

\subsection{什么是Pandas？}

\begin{knowledgebox}[知识点：Pandas是什么]
Pandas是Python的一个数据分析库，它的名字来源于"Panel Data"（面板数据）。

简单来说，Pandas就是用来处理表格数据的工具，就像Excel一样，但比Excel更强大、更灵活。

\textbf{为什么叫Pandas？}
\begin{itemize}
    \item 和熊猫（panda）没有关系
    \item 来自计量经济学的"Panel Data"概念
    \item 是Python Data Analysis的缩写
\end{itemize}
\end{knowledgebox}

\subsection{安装Pandas}

在使用Pandas之前，需要先安装它。打开命令行（CMD或PowerShell），输入以下命令：

\begin{lstlisting}[caption={安装Pandas}]
# 使用pip安装（推荐）
pip install pandas

# 如果安装了Anaconda，使用conda安装
conda install pandas
\end{lstlisting}

\begin{tipbox}[安装小贴士]
\begin{itemize}
    \item 安装时需要联网
    \item 如果安装速度慢，可以使用国内镜像源
    \item 安装完成后，关闭命令行重新打开
\end{itemize}
\end{tipbox}

\subsection{导入Pandas}

安装完成后，在Python代码中导入Pandas：

\begin{lstlisting}[caption={导入Pandas}]
# 导入pandas库，并给它起别名pd
import pandas as pd

# 同时导入numpy（经常和pandas一起使用）
import numpy as np
\end{lstlisting}

\begin{knowledgebox}[知识点：为什么要用别名？]
在Python中，给库起别名是为了：
\begin{enumerate}
    \item \textbf{简化代码}：写\texttt{pd}比写\texttt{pandas}更短
    \item \textbf{行业惯例}：大家都用\texttt{pd}，方便阅读
    \item \textbf{保持一致}：不同人写的代码风格统一
\end{enumerate}

这就好比：
\begin{itemize}
    \item 你的名字叫"张三丰"，但朋友们叫你"小张"
    \item 写代码时，\texttt{pandas}是全名，\texttt{pd}是小名
\end{itemize}
\end{knowledgebox}

%========================================
\section{Pandas的核心概念}
%========================================

\subsection{DataFrame：数据表格}

\begin{knowledgebox}[知识点：什么是DataFrame？]
DataFrame是Pandas中最重要的数据结构，你可以把它想象成一张Excel表格：

\begin{itemize}
    \item 有行和列
    \item 每列可以是不同的数据类型（数字、文字、日期等）
    \item 有列名（表头）
    \item 每行有一个索引（行号）
\end{itemize}

\textbf{类比理解：}
\begin{itemize}
    \item Excel表格 $\rightarrow$ DataFrame
    \item 工作表 $\rightarrow$ DataFrame
    \item 数据库表 $\rightarrow$ DataFrame
\end{itemize}
\end{knowledgebox}

\subsection{创建第一个DataFrame}

让我们创建一个简单的学生成绩表：

\begin{lstlisting}[caption={创建DataFrame}]
import pandas as pd

# 创建一个字典（dict），存储数据
data = {
    '姓名': ['张三', '李四', '王五'],
    '数学': [85, 92, 78],
    '英语': [88, 85, 92]
}

# 用字典创建DataFrame
df = pd.DataFrame(data)

# 显示DataFrame
print(df)
\end{lstlisting}

运行结果：
\begin{verbatim}
   姓名  数学  英语
0  张三   85   88
1  李四   92   85
2  王五   78   92
\end{verbatim}

\begin{knowledgebox}[知识点：代码拆解]
让我们逐行理解上面的代码：

\begin{enumerate}
    \item \texttt{data = \{...\}}：创建一个字典
    \begin{itemize}
        \item 字典的键（'姓名'、'数学'）变成列名
        \item 字典的值（列表）变成列的数据
    \end{itemize}
    
    \item \texttt{pd.DataFrame(data)}：把字典转换成DataFrame
    
    \item \texttt{print(df)}：显示表格内容
    \begin{itemize}
        \item 最左边一列0、1、2是\textbf{行索引}（默认从0开始）
        \item 上面一行是\textbf{列名}
    \end{itemize}
\end{enumerate}
\end{knowledgebox}

\begin{practicebox}[动手练习1]
试着创建一个包含以下信息的DataFrame：
\begin{itemize}
    \item 3个商品：手机、电脑、耳机
    \item 每个商品有：名称、价格、库存数量
\end{itemize}

\textbf{参考代码：}
\begin{lstlisting}
import pandas as pd

data = {
    '商品名称': ['手机', '电脑', '耳机'],
    '价格': [2999, 5999, 199],
    '库存': [100, 50, 200]
}

df = pd.DataFrame(data)
print(df)
\end{lstlisting}
\end{practicebox}

%========================================
\section{读取数据文件}
%========================================

\subsection{为什么要从文件读取数据？}

在实际工作中，数据通常存储在文件中，而不是直接写在代码里。常见的数据文件格式有：

\begin{itemize}
    \item \textbf{CSV文件}（.csv）：逗号分隔的值，最常用
    \item \textbf{Excel文件}（.xlsx）：Excel表格
    \item \textbf{JSON文件}（.json）：网页数据常用格式
\end{itemize}

\subsection{读取CSV文件}

CSV（Comma-Separated Values）是最常见的数据格式。假设我们有一个文件\texttt{学生成绩.csv}：

\begin{lstlisting}[caption={读取CSV文件}]
import pandas as pd

# 读取CSV文件
df = pd.read_csv('学生成绩.csv')

# 显示前5行
print(df.head())
\end{lstlisting}

\begin{knowledgebox}[知识点：read\_csv常用参数]
\texttt{pd.read\_csv()}有很多参数，常用的有：

\begin{itemize}
    \item \texttt{encoding='utf-8'}：指定文件编码（处理中文）
    \item \texttt{sep=','}：指定分隔符（默认是逗号）
    \item \texttt{header=0}：指定第几行是列名（默认第0行）
\end{itemize}

\textbf{示例：}
\begin{lstlisting}
# 读取中文CSV文件（防止乱码）
df = pd.read_csv('学生成绩.csv', encoding='utf-8')

# 读取用制表符分隔的文件
df = pd.read_csv('数据.txt', sep='\t')
\end{lstlisting}
\end{knowledgebox}

\begin{tipbox}[中文乱码怎么办？]
如果读取后中文显示乱码，尝试以下编码：
\begin{itemize}
    \item \texttt{encoding='utf-8'}：最常用的编码
    \item \texttt{encoding='gbk'}：中文Windows系统常用
    \item \texttt{encoding='gb2312'}：简体中文编码
\end{itemize}
\end{tipbox}

\subsection{查看数据基本信息}

读取数据后，首先要了解数据的基本情况：

\begin{lstlisting}[caption={查看数据信息}]
# 查看数据形状（行数，列数）
print(df.shape)

# 查看列名
print(df.columns)

# 查看数据类型
print(df.dtypes)

# 查看基本信息（包括非空值数量）
print(df.info())

# 查看统计摘要（只针对数值列）
print(df.describe())
\end{lstlisting}

\begin{knowledgebox}[知识点：这些方法的作用]
\begin{itemize}
    \item \texttt{shape}：返回\texttt{(行数, 列数)}，比如\texttt{(100, 5)}表示100行5列
    \item \texttt{columns}：返回所有列名
    \item \texttt{dtypes}：返回每列的数据类型
    \item \texttt{info()}：显示详细信息，包括每列的非空值数量
    \item \texttt{describe()}：统计数值列的计数、均值、最值等
\end{itemize}
\end{knowledgebox}

\begin{practicebox}[动手练习2]
使用提供的\texttt{01\_电商订单数据\_练习.csv}文件：

\begin{enumerate}
    \item 读取这个CSV文件
    \item 查看数据有多少行、多少列
    \item 查看所有列名
    \item 查看每列的数据类型
\end{enumerate}

\textbf{参考代码：}
\begin{lstlisting}
import pandas as pd

# 读取数据
df = pd.read_csv('01_电商订单数据_练习.csv')

# 查看形状
print("数据形状:", df.shape)

# 查看列名
print("\n列名:")
print(df.columns)

# 查看数据类型
print("\n数据类型:")
print(df.dtypes)
\end{lstlisting}
\end{practicebox}

%========================================
\section{数据选择：如何选取需要的数据}
%========================================

\subsection{选择列}

有时候我们只需要表格中的某些列，而不是全部。

\subsubsection{选择单列}

\begin{lstlisting}[caption={选择单列}]
# 选择'姓名'这一列
names = df['姓名']

# 或者使用点号（仅适用于列名没有空格和特殊字符时）
names = df.姓名
\end{lstlisting}

\begin{knowledgebox}[知识点：单列选择的结果]
选择单列后，得到的是一个\textbf{Series}（序列），不是DataFrame。

\textbf{Series和DataFrame的区别：}
\begin{itemize}
    \item DataFrame：二维表格（有行有列）
    \item Series：一维数组（只有一列数据）
\end{itemize}
\end{knowledgebox}

\subsubsection{选择多列}

\begin{lstlisting}[caption={选择多列}]
# 选择多列（用两个方括号）
subset = df[['姓名', '数学']]

# 注意：外层[]是DataFrame的索引，内层[]是Python列表
\end{lstlisting}

\begin{tipbox}[注意双括号]
选择多列时，需要用\textbf{两个方括号}：
\begin{itemize}
    \item 外层\texttt{[]}：DataFrame的列选择
    \item 内层\texttt{[]}：Python的列表，包含要选的列名
\end{itemize}

\textbf{常见错误：}
\begin{lstlisting}
# 错误！只用一个方括号
df['姓名', '数学']  # 会报错

# 正确！用两个方括号
df[['姓名', '数学']]  # OK
\end{lstlisting}
\end{tipbox}

\subsection{选择行}

Pandas提供了两种选择行的方法：

\begin{knowledgebox}[知识点：loc vs iloc]
\begin{itemize}
    \item \texttt{loc}：按\textbf{标签}选择（用行名）
    \item \texttt{iloc}：按\textbf{位置}选择（用行号，从0开始）
\end{itemize}

\textbf{类比：}
\begin{itemize}
    \item \texttt{loc}就像用姓名找人
    \item \texttt{iloc}就像用学号找人（第1个、第2个...）
\end{itemize}
\end{knowledgebox}

\subsubsection{使用iloc选择行}

\begin{lstlisting}[caption={使用iloc选择行}]
# 选择第1行（索引为0）
row = df.iloc[0]

# 选择前3行
rows = df.iloc[0:3]

# 选择第1、3、5行
rows = df.iloc[[0, 2, 4]]
\end{lstlisting}

\subsubsection{使用loc选择行}

\begin{lstlisting}[caption={使用loc选择行}]
# 假设行索引是姓名
df_indexed = df.set_index('姓名')

# 选择姓名为'张三'的行
row = df_indexed.loc['张三']

# 选择多个人的数据
rows = df_indexed.loc[['张三', '李四']]
\end{lstlisting}

\subsection{同时选择行和列}

\begin{lstlisting}[caption={同时选择行和列}]
# iloc[行位置, 列位置]
# 选择前3行，前2列
data = df.iloc[0:3, 0:2]

# loc[行标签, 列标签]
# 选择'张三'和'李四'的'数学'和'英语'成绩
data = df_indexed.loc[['张三', '李四'], ['数学', '英语']]
\end{lstlisting}

\begin{practicebox}[动手练习3]
使用\texttt{01\_电商订单数据\_练习.csv}：

\begin{enumerate}
    \item 只选择\texttt{product\_name}和\texttt{price}两列
    \item 选择前10行数据
    \item 选择第5到第15行的\texttt{order\_id}、\texttt{user\_id}、\texttt{quantity}三列
\end{enumerate}

\textbf{参考代码：}
\begin{lstlisting}
import pandas as pd

df = pd.read_csv('01_电商订单数据_练习.csv')

# 1. 选择两列
product_info = df[['product_name', 'price']]
print(product_info.head())

# 2. 选择前10行
first_10 = df.iloc[0:10]
print(first_10)

# 3. 选择特定行和列
subset = df.iloc[4:15][['order_id', 'user_id', 'quantity']]
print(subset)
\end{lstlisting}
\end{practicebox}

%========================================
\section{数据清洗：处理脏数据}
%========================================

\subsection{什么是脏数据？}

\begin{knowledgebox}[知识点：常见的数据质量问题]
真实世界的数据往往存在各种问题，我们称之为"脏数据"：

\begin{enumerate}
    \item \textbf{缺失值}：某些单元格是空的
    \item \textbf{重复值}：同样的数据出现多次
    \item \textbf{异常值}：明显不合理的数据（如年龄为-5）
    \item \textbf{格式不一致}：日期格式混乱、大小写不统一
    \item \textbf{类型错误}：数字被存成了文字
\end{enumerate}
\end{knowledgebox}

\subsection{处理缺失值}

\subsubsection{识别缺失值}

\begin{lstlisting}[caption={检查缺失值}]
# 检查每列的缺失值数量
print(df.isnull().sum())

# 检查缺失值比例
missing_ratio = df.isnull().sum() / len(df) * 100
print(missing_ratio)
\end{lstlisting}

\begin{knowledgebox}[知识点：isnull()的作用]
\texttt{isnull()}方法会返回一个和原表格一样大的表格，里面只有True和False：
\begin{itemize}
    \item True：表示这个单元格是缺失值
    \item False：表示这个单元格有数据
\end{itemize}

\texttt{sum()}会把True当成1、False当成0，所以能统计出每列的缺失值数量。
\end{knowledgebox}

\subsubsection{填充缺失值}

有几种常见的填充策略：

\begin{lstlisting}[caption={填充缺失值的方法}]
# 方法1：用固定值填充
df['年龄'] = df['年龄'].fillna(0)

# 方法2：用平均值填充（适合数值列）
df['成绩'] = df['成绩'].fillna(df['成绩'].mean())

# 方法3：用中位数填充（适合有异常值的情况）
df['工资'] = df['工资'].fillna(df['工资'].median())

# 方法4：用众数填充（适合类别数据）
df['性别'] = df['性别'].fillna(df['性别'].mode()[0])

# 方法5：用前一个值填充（适合时间序列）
df['温度'] = df['温度'].fillna(method='ffill')
\end{lstlisting}

\begin{tipbox}[选择填充策略的建议]
\begin{itemize}
    \item \textbf{数值数据}：优先用中位数（对异常值不敏感）
    \item \textbf{类别数据}：用众数（出现次数最多的值）
    \item \textbf{时间序列}：用前后值填充（保持趋势）
    \item \textbf{缺失太多}（如超过50\%）：考虑删除该列
\end{itemize}
\end{tipbox}

\subsubsection{删除缺失值}

如果缺失值很少，可以直接删除：

\begin{lstlisting}[caption={删除缺失值}]
# 删除包含缺失值的行
df_clean = df.dropna()

# 只删除某列有缺失值的行
df_clean = df.dropna(subset=['年龄'])

# 删除整列都是缺失值的列
df_clean = df.dropna(axis=1, how='all')
\end{lstlisting}

\begin{practicebox}[动手练习4]
使用\texttt{03\_缺失值处理\_练习.csv}：

\begin{enumerate}
    \item 统计每列的缺失值数量
    \item 用中位数填充\texttt{salary}列的缺失值
    \item 用众数填充\texttt{department}列的缺失值
    \item 删除\texttt{score}列仍有缺失值的行
\end{enumerate}

\textbf{参考代码：}
\begin{lstlisting}
import pandas as pd

df = pd.read_csv('03_缺失值处理_练习.csv')

# 1. 统计缺失值
print("缺失值统计:")
print(df.isnull().sum())

# 2. 用中位数填充工资
median_salary = df['salary'].median()
df['salary'] = df['salary'].fillna(median_salary)
print(f"\n用 {median_salary} 填充工资缺失值")

# 3. 用众数填充部门
mode_dept = df['department'].mode()[0]
df['department'] = df['department'].fillna(mode_dept)
print(f"用 '{mode_dept}' 填充部门缺失值")

# 4. 删除score仍有缺失的行
df = df.dropna(subset=['score'])
print(f"\n删除后剩余 {len(df)} 行")
\end{lstlisting}
\end{practicebox}

\subsection{处理重复值}

\begin{lstlisting}[caption={处理重复值}]
# 检查重复值
duplicates = df.duplicated()
print(f"有 {duplicates.sum()} 行重复")

# 查看重复的行
duplicate_rows = df[df.duplicated(keep=False)]
print(duplicate_rows)

# 删除重复值（保留第一个）
df_clean = df.drop_duplicates()

# 基于特定列删除重复
df_clean = df.drop_duplicates(subset=['姓名'])
\end{lstlisting}

\subsection{处理异常值}

\subsubsection{什么是异常值？}

异常值是指明显偏离正常范围的数据，比如：
\begin{itemize}
    \item 年龄为负数或超过150
    \item 价格为0或异常高
    \item 数量为负数
\end{itemize}

\subsubsection{检测异常值}

\begin{lstlisting}[caption={使用IQR方法检测异常值}]
# IQR（四分位距）方法
def detect_outliers_iqr(data):
    Q1 = data.quantile(0.25)  # 第25百分位数
    Q3 = data.quantile(0.75)  # 第75百分位数
    IQR = Q3 - Q1
    
    # 异常值范围
    lower_bound = Q1 - 1.5 * IQR
    upper_bound = Q3 + 1.5 * IQR
    
    # 找出异常值
    outliers = data[(data < lower_bound) | (data > upper_bound)]
    return outliers

# 检测价格的异常值
price_outliers = detect_outliers_iqr(df['price'])
print(f"发现 {len(price_outliers)} 个价格异常值")
\end{lstlisting}

\begin{knowledgebox}[知识点：IQR方法原理]
IQR（Interquartile Range，四分位距）是一种稳健的异常值检测方法：

\begin{enumerate}
    \item 把数据从小到大排序
    \item 找到第25\%位置的数（Q1）和第75\%位置的数（Q3）
    \item 计算IQR = Q3 - Q1
    \item 正常数据范围：[Q1 - 1.5×IQR, Q3 + 1.5×IQR]
    \item 超出这个范围的就是异常值
\end{enumerate}

\textbf{优点}：不受极端值影响，适合偏态分布的数据
\end{knowledgebox}

\subsubsection{处理异常值}

找到异常值后，有几种处理方式：

\begin{lstlisting}[caption={处理异常值的方法}]
# 方法1：删除异常值
df_clean = df[(df['price'] >= lower_bound) & (df['price'] <= upper_bound)]

# 方法2：用边界值替换（截断）
df['price'] = df['price'].clip(lower=lower_bound, upper=upper_bound)

# 方法3：用中位数替换
median_price = df['price'].median()
df.loc[df['price'] > upper_bound, 'price'] = median_price
\end{lstlisting}

%========================================
\section{数据转换与格式化}
%========================================

\subsection{数据类型转换}

\begin{lstlisting}[caption={转换数据类型}]
# 查看当前数据类型
print(df.dtypes)

# 转换为整数
df['年龄'] = df['年龄'].astype(int)

# 转换为浮点数
df['价格'] = df['价格'].astype(float)

# 转换为字符串
df['编号'] = df['编号'].astype(str)

# 转换为日期时间
df['日期'] = pd.to_datetime(df['日期'])
\end{lstlisting}

\subsection{字符串处理}

\begin{lstlisting}[caption={常用字符串方法}]
# 去除空格
df['姓名'] = df['姓名'].str.strip()

# 转换大小写
df['邮箱'] = df['邮箱'].str.lower()

# 替换字符
df['电话'] = df['电话'].str.replace('-', '')

# 提取子串
df['年份'] = df['日期'].str[0:4]

# 分割字符串
df['姓名'].str.split(' ', expand=True)
\end{lstlisting}

\begin{knowledgebox}[知识点：str访问器]
在Pandas中，要对字符串列进行操作，需要使用\texttt{.str}访问器：

\begin{itemize}
    \item \texttt{df['列名'].str.方法()}
    \item 不能直接写\texttt{df['列名'].strip()}
\end{itemize}

这就像打开一个工具箱，\texttt{.str}就是打开字符串工具箱的钥匙。
\end{knowledgebox}

\subsection{应用函数}

有时候需要对数据进行复杂的转换，可以自定义函数：

\begin{lstlisting}[caption={使用apply应用函数}]
# 定义一个函数
def score_level(score):
    if score >= 90:
        return '优秀'
    elif score >= 80:
        return '良好'
    elif score >= 60:
        return '及格'
    else:
        return '不及格'

# 应用到整列
df['等级'] = df['成绩'].apply(score_level)

# 使用lambda函数（简单操作）
df['价格_打折'] = df['价格'].apply(lambda x: x * 0.8)
\end{lstlisting}

\begin{practicebox}[动手练习5]
使用\texttt{04\_数据格式清洗\_练习.csv}：

\begin{enumerate}
    \item 去除\texttt{name}列的前后空格
    \item 将\texttt{income}列中的"k"替换为"000"，并转换为数值
    \item 从\texttt{birth\_date}列中提取年份
\end{enumerate}

\textbf{参考代码：}
\begin{lstlisting}
import pandas as pd

df = pd.read_csv('04_数据格式清洗_练习.csv')

# 1. 去除空格
df['name'] = df['name'].str.strip()

# 2. 处理收入（将12k转换为12000）
def convert_income(income_str):
    income_str = str(income_str)
    if 'k' in income_str.lower():
        number = float(income_str.lower().replace('k', ''))
        return number * 1000
    else:
        return float(income_str)

df['income_num'] = df['income'].apply(convert_income)

# 3. 提取年份
df['birth_year'] = df['birth_date'].str[0:4]

print(df[['name', 'income', 'income_num', 'birth_year']].head())
\end{lstlisting}
\end{practicebox}

%========================================
\section{数据分组与统计}
%========================================

\subsection{分组统计}

\begin{knowledgebox}[知识点：分组统计的概念]
分组统计就是：
\begin{enumerate}
    \item 先把数据按某个列分成若干组
    \item 然后对每组进行统计计算
\end{enumerate}

\textbf{类比}：
\begin{itemize}
    \item 把全班同学按\textbf{性别}分组，计算每组平均成绩
    \item 把商品按\textbf{类别}分组，计算每类总销量
\end{itemize}
\end{knowledgebox}

\begin{lstlisting}[caption={分组统计}]
# 按部门分组，计算平均工资
dept_avg = df.groupby('部门')['工资'].mean()
print(dept_avg)

# 多列分组，多统计量
df.groupby(['部门', '性别']).agg({
    '工资': ['mean', 'max', 'min'],
    '年龄': 'mean'
})
\end{lstlisting}

\begin{knowledgebox}[知识点：groupby的工作流程]
\texttt{groupby()}的工作流程就像流水线：

\begin{enumerate}
    \item \textbf{分组}：按指定列的值把数据分成多组
    \item \textbf{应用}：对每组应用统计函数
    \item \textbf{合并}：把各组结果合并成一个表格
\end{enumerate}

\textbf{常用统计函数：}
\begin{itemize}
    \item \texttt{mean()}：平均值
    \item \texttt{sum()}：求和
    \item \texttt{count()}：计数
    \item \texttt{max()/min()}：最大/最小值
    \item \texttt{std()}：标准差
\end{itemize}
\end{knowledgebox}

%========================================
\section{综合实战：完整的数据处理流程}
%========================================

\subsection{案例背景}

假设我们有一个电商订单数据集，需要完成以下任务：
\begin{enumerate}
    \item 清洗数据（处理缺失值、重复值、异常值）
    \item 转换数据格式（价格、日期、手机号）
    \item 分析数据（统计各品类销售额）
\end{enumerate}

\subsection{分步实现}

\begin{lstlisting}[caption={完整数据处理流程}]
import pandas as pd
import numpy as np

# ========== 第1步：读取数据 ==========
print("步骤1：读取数据")
df = pd.read_csv('01_电商订单数据_练习.csv')
print(f"原始数据：{df.shape[0]}行，{df.shape[1]}列")

# ========== 第2步：查看数据质量 ==========
print("\n步骤2：数据质量检查")
print("缺失值统计：")
print(df.isnull().sum())
print(f"\n重复值：{df.duplicated().sum()}行")

# ========== 第3步：处理重复值 ==========
print("\n步骤3：删除重复值")
df = df.drop_duplicates()
print(f"删除后：{df.shape[0]}行")

# ========== 第4步：处理缺失值 ==========
print("\n步骤4：处理缺失值")
# user_id缺失用'UNKNOWN'填充
df['user_id'] = df['user_id'].fillna('UNKNOWN')

# ========== 第5步：处理异常值 ==========
print("\n步骤5：处理异常值")
# 处理负数数量
df.loc[df['quantity'] < 0, 'quantity'] = 1
print("负数数量已修正为1")

# ========== 第6步：格式转换 ==========
print("\n步骤6：格式转换")
# 提取价格数字
df['price_num'] = df['price'].str.replace('元', '').astype(float)

# 标准化手机号
df['phone_clean'] = df['phone'].str.replace('-', '')

# ========== 第7步：数据分析 ==========
print("\n步骤7：数据分析")
# 计算订单金额
df['total_amount'] = df['price_num'] * df['quantity']

# 按商品统计
product_stats = df.groupby('product_name').agg({
    'total_amount': 'sum',
    'quantity': 'sum'
}).sort_values('total_amount', ascending=False)

print("\n销售额TOP5商品：")
print(product_stats.head())

# ========== 第8步：保存结果 ==========
print("\n步骤8：保存结果")
df.to_csv('清洗后数据.csv', index=False)
print("数据已保存到：清洗后数据.csv")
\end{lstlisting}

%========================================
\section{竞赛模拟题}
%========================================

\subsection{题目一：电商订单数据处理}

\begin{practicebox}[竞赛模拟题1]
\textbf{题目要求：}

使用\texttt{01\_电商订单数据\_练习.csv}，完成以下任务：

\begin{enumerate}
    \item 读取数据，查看基本信息
    \item 删除重复订单
    \item 处理user\_id缺失值（用'UNKNOWN'填充）
    \item 处理负数数量（修正为1）
    \item 提取价格中的数字，转换为数值型
    \item 标准化手机号格式（去除横线）
    \item 计算每个订单的总金额（价格×数量）
    \item 统计每个商品的总销售额
    \item 保存清洗后的数据
\end{enumerate}

\textbf{输出要求：}
\begin{itemize}
    \item 清洗后的数据文件
    \item 数据统计报告（包含处理记录和统计结果）
\end{itemize}
\end{practicebox}

\subsection{题目二：员工信息数据处理}

\begin{practicebox}[竞赛模拟题2]
\textbf{题目要求：}

使用\texttt{03\_缺失值处理\_练习.csv}和\texttt{04\_数据格式清洗\_练习.csv}，完成以下任务：

\begin{enumerate}
    \item 读取两个CSV文件
    \item 处理缺失值：
    \begin{itemize}
        \item age：用中位数填充
        \item salary：用平均值填充
        \item department：用众数填充
    \end{itemize}
    \item 格式标准化：
    \begin{itemize}
        \item 姓名：去除前后空格
        \item 收入：统一转换为数值（处理"k"单位）
    \end{itemize}
    \item 按部门统计平均工资和年龄
    \item 保存处理后的数据
\end{enumerate}
\end{practicebox}

\subsection{题目一参考答案}

\begin{practicebox}[竞赛模拟题1 - 完整解答]

\textbf{解题思路：}

这道题要求我们对电商订单数据进行全面清洗和分析。按照数据处理的正常流程，我们应该：
\begin{enumerate}
    \item 先读取数据，了解数据的基本情况
    \item 处理数据质量问题（重复、缺失、异常）
    \item 进行格式转换
    \item 进行数据分析
    \item 保存结果
\end{enumerate}

\textbf{完整代码：}

\begin{lstlisting}[caption={电商订单数据处理完整解答}]
import pandas as pd
import numpy as np

# ========== 第1步：读取数据 ==========
print("=" * 50)
print("步骤1：读取数据")
print("=" * 50)

df = pd.read_csv('01_电商订单数据_练习.csv')
print(f"数据形状：{df.shape}")
print(f"\n列名：{df.columns.tolist()}")
print("\n前5行数据：")
print(df.head())

# ========== 第2步：删除重复值 ==========
print("\n" + "=" * 50)
print("步骤2：删除重复值")
print("=" * 50)

# 查看重复值
n_duplicates = df.duplicated().sum()
print(f"发现 {n_duplicates} 行重复数据")

# 删除重复值
df = df.drop_duplicates()
print(f"删除后剩余 {len(df)} 行")

# ========== 第3步：处理缺失值 ==========
print("\n" + "=" * 50)
print("步骤3：处理缺失值")
print("=" * 50)

# 查看缺失值
print("缺失值统计：")
print(df.isnull().sum())

# 填充user_id缺失值
df['user_id'] = df['user_id'].fillna('UNKNOWN')
print("\nuser_id缺失值已用'UNKNOWN'填充")

# ========== 第4步：处理异常值 ==========
print("\n" + "=" * 50)
print("步骤4：处理异常值")
print("=" * 50)

# 查看数量列的情况
print("数量统计：")
print(df['quantity'].describe())

# 处理负数数量
negative_count = (df['quantity'] < 0).sum()
print(f"\n发现 {negative_count} 个负数数量")

# 将负数改为1
df.loc[df['quantity'] < 0, 'quantity'] = 1
print("负数数量已修正为1")

# ========== 第5步：提取价格数字 ==========
print("\n" + "=" * 50)
print("步骤5：提取价格数字")
print("=" * 50)

# 查看价格列的格式
print("价格示例：")
print(df['price'].head())

# 提取数字（去除'元'字）
df['price_num'] = df['price'].str.replace('元', '').astype(float)

print("\n转换后的价格：")
print(df[['price', 'price_num']].head())

# ========== 第6步：标准化手机号 ==========
print("\n" + "=" * 50)
print("步骤6：标准化手机号")
print("=" * 50)

# 查看手机号格式
print("手机号示例：")
print(df['phone'].head())

# 去除横线
df['phone_clean'] = df['phone'].str.replace('-', '')

print("\n清洗后的手机号：")
print(df[['phone', 'phone_clean']].head())

# ========== 第7步：计算订单金额 ==========
print("\n" + "=" * 50)
print("步骤7：计算订单金额")
print("=" * 50)

# 计算总金额 = 价格 × 数量
df['total_amount'] = df['price_num'] * df['quantity']

print("订单金额计算完成")
print(df[['product_name', 'price_num', 'quantity', 'total_amount']].head())

# ========== 第8步：统计商品销售额 ==========
print("\n" + "=" * 50)
print("步骤8：统计商品销售额")
print("=" * 50)

# 按商品分组统计
product_sales = df.groupby('product_name').agg({
    'total_amount': 'sum',      # 总销售额
    'quantity': 'sum',          # 总销量
    'order_id': 'count'         # 订单数
}).sort_values('total_amount', ascending=False)

product_sales.columns = ['总销售额', '总销量', '订单数']

print("商品销售额TOP10：")
print(product_sales.head(10))

# ========== 第9步：生成报告并保存 ==========
print("\n" + "=" * 50)
print("步骤9：生成报告并保存")
print("=" * 50)

# 生成处理报告
report = f"""
数据清洗报告
{'=' * 50}
1. 原始数据：50行，9列
2. 删除重复值：{n_duplicates}行
3. 处理后数据：{len(df)}行
4. 处理内容：
   - user_id缺失值用'UNKNOWN'填充
   - 负数数量修正为1
   - 价格转换为数值型
   - 手机号去除横线
5. 统计结果：
   - 商品种类：{df['product_name'].nunique()}种
   - 总销售额：{df['total_amount'].sum():.2f}元
   - 总订单数：{len(df)}单
{'=' * 50}
"""

print(report)

# 保存报告到文件
with open('数据处理报告.txt', 'w', encoding='utf-8') as f:
    f.write(report)

print("报告已保存到：数据处理报告.txt")

# 保存清洗后的数据
df.to_csv('清洗后订单数据.csv', index=False, encoding='utf-8')
print("数据已保存到：清洗后订单数据.csv")

print("\n" + "=" * 50)
print("数据处理完成！")
print("=" * 50)
\end{lstlisting}

\textbf{关键步骤说明：}

\begin{enumerate}
    \item \textbf{读取数据}：使用\texttt{pd.read\_csv()}读取CSV文件
    \item \textbf{查看信息}：用\texttt{shape}、\texttt{columns}、\texttt{head()}了解数据
    \item \textbf{删除重复}：\texttt{duplicated()}检测，\texttt{drop\_duplicates()}删除
    \item \textbf{填充缺失}：\texttt{fillna()}用指定值填充
    \item \textbf{处理异常}：用条件筛选找到异常值，然后修正
    \item \textbf{字符串处理}：\texttt{str.replace()}去除不需要的字符
    \item \textbf{类型转换}：\texttt{astype(float)}转换为数值型
    \item \textbf{计算新列}：直接用列相乘生成新列
    \item \textbf{分组统计}：\texttt{groupby()}配合\texttt{agg()}进行多维度统计
\end{enumerate}
\end{practicebox}

\subsection{题目二参考答案}

\begin{practicebox}[竞赛模拟题2 - 完整解答]

\textbf{解题思路：}

这道题涉及两个文件的处理，需要：
\begin{enumerate}
    \item 分别读取两个CSV文件
    \item 对缺失值采用不同的填充策略
    \item 对格式混乱的数据进行标准化
    \item 合并分析结果
\end{enumerate}

\textbf{完整代码：}

\begin{lstlisting}[caption={员工信息数据处理完整解答}]
import pandas as pd
import numpy as np

# ========== 第1步：读取数据 ==========
print("=" * 50)
print("步骤1：读取数据")
print("=" * 50)

# 读取两个CSV文件
df1 = pd.read_csv('03_缺失值处理_练习.csv')
df2 = pd.read_csv('04_数据格式清洗_练习.csv')

print(f"员工数据：{df1.shape}")
print(f"格式清洗数据：{df2.shape}")

# ========== 第2步：处理缺失值（df1） ==========
print("\n" + "=" * 50)
print("步骤2：处理缺失值")
print("=" * 50)

# 查看缺失值情况
print("缺失值统计：")
print(df1.isnull().sum())

# 2.1 age用中位数填充
age_median = df1['age'].median()
df1['age'] = df1['age'].fillna(age_median)
print(f"\nage缺失值用中位数 {age_median} 填充")

# 2.2 salary用平均值填充
salary_mean = df1['salary'].mean()
df1['salary'] = df1['salary'].fillna(salary_mean)
print(f"salary缺失值用平均值 {salary_mean:.2f} 填充")

# 2.3 department用众数填充
dept_mode = df1['department'].mode()[0]
df1['department'] = df1['department'].fillna(dept_mode)
print(f"department缺失值用众数 '{dept_mode}' 填充")

# 验证缺失值是否处理完毕
print("\n处理后缺失值统计：")
print(df1.isnull().sum())

# ========== 第3步：格式标准化（df2） ==========
print("\n" + "=" * 50)
print("步骤3：格式标准化")
print("=" * 50)

# 3.1 去除姓名前后空格
print("姓名处理前示例：")
print(df2['name'].head())

df2['name'] = df2['name'].str.strip()

print("\n姓名处理后示例：")
print(df2['name'].head())

# 3.2 处理收入格式（将12k转换为12000）
print("\n收入处理前示例：")
print(df2['income'].head(10))

def convert_income(income_str):
    """将收入转换为数值"""
    # 转换为字符串
    income_str = str(income_str).strip()
    
    # 如果包含k或K，表示千元
    if 'k' in income_str.lower():
        # 提取数字部分
        number = float(income_str.lower().replace('k', ''))
        return number * 1000
    else:
        # 直接转换为浮点数
        return float(income_str)

# 应用转换函数
df2['income_num'] = df2['income'].apply(convert_income)

print("\n收入处理后示例：")
print(df2[['income', 'income_num']].head(10))

# ========== 第4步：按部门统计 ==========
print("\n" + "=" * 50)
print("步骤4：按部门统计")
print("=" * 50)

# 4.1 员工数据统计
dept_stats = df1.groupby('department').agg({
    'salary': ['mean', 'max', 'min'],
    'age': 'mean'
}).round(2)

# 重命名列
dept_stats.columns = ['平均工资', '最高工资', '最低工资', '平均年龄']

print("各部门统计：")
print(dept_stats)

# 4.2 计算部门人数
dept_count = df1['department'].value_counts()
print("\n部门人数：")
print(dept_count)

# ========== 第5步：生成完整报告 ==========
print("\n" + "=" * 50)
print("步骤5：生成报告")
print("=" * 50)

report = f"""
员工数据处理报告
{'=' * 50}

一、数据概况
  员工数据：{len(df1)}人
  部门数量：{df1['department'].nunique()}个

二、缺失值处理
  age：用中位数 {age_median} 填充
  salary：用平均值 {salary_mean:.2f} 填充
  department：用众数 '{dept_mode}' 填充

三、格式标准化
  姓名：去除前后空格
  收入：统一转换为数值（处理k单位）

四、部门统计
{dept_stats.to_string()}

五、部门人数
{dept_count.to_string()}
{'=' * 50}
"""

print(report)

# 保存报告
with open('员工数据处理报告.txt', 'w', encoding='utf-8') as f:
    f.write(report)

print("报告已保存到：员工数据处理报告.txt")

# ========== 第6步：保存处理后的数据 ==========
print("\n" + "=" * 50)
print("步骤6：保存数据")
print("=" * 50)

# 保存员工数据
df1.to_csv('处理后员工数据.csv', index=False, encoding='utf-8')
print("员工数据已保存到：处理后员工数据.csv")

# 保存格式清洗数据
df2.to_csv('处理后格式数据.csv', index=False, encoding='utf-8')
print("格式数据已保存到：处理后格式数据.csv")

print("\n" + "=" * 50)
print("数据处理完成！")
print("=" * 50)
\end{lstlisting}

\textbf{关键步骤说明：}

\begin{enumerate}
    \item \textbf{读取多个文件}：分别用\texttt{pd.read\_csv()}读取
    \item \textbf{不同填充策略}：
    \begin{itemize}
        \item 数值型（age）：用\texttt{median()}中位数，不受异常值影响
        \item 数值型（salary）：用\texttt{mean()}平均值
        \item 分类型（department）：用\texttt{mode()}众数
    \end{itemize}
    \item \textbf{自定义函数}：用\texttt{apply()}应用自定义转换函数
    \item \textbf{多维度统计}：用\texttt{agg()}同时进行多个统计计算
    \item \textbf{生成报告}：将处理过程和结果整理成文本报告
\end{enumerate}
\end{practicebox}

%========================================
\section{常用代码速查表}
%========================================

\subsection{数据读取与保存}

\begin{lstlisting}[caption={读写数据}]
# 读取CSV
df = pd.read_csv('文件.csv', encoding='utf-8')

# 保存CSV
df.to_csv('输出.csv', index=False, encoding='utf-8')
\end{lstlisting}

\subsection{数据查看}

\begin{lstlisting}[caption={查看数据}]
df.head()           # 前5行
df.tail()           # 后5行
df.shape            # 形状（行数，列数）
df.columns          # 列名
df.dtypes           # 数据类型
df.info()           # 详细信息
df.describe()       # 统计摘要
\end{lstlisting}

\subsection{数据选择}

\begin{lstlisting}[caption={选择数据}]
# 选择列
df['列名']          # 单列
df[['列1', '列2']]  # 多列

# 选择行
df.iloc[0]          # 第1行（按位置）
df.iloc[0:5]        # 前5行
df.loc['标签']      # 按标签

# 条件筛选
df[df['年龄'] > 18]
df[(df['年龄'] > 18) & (df['性别'] == '男')]
\end{lstlisting}

\subsection{数据清洗}

\begin{lstlisting}[caption={清洗数据}]
# 缺失值
df.isnull().sum()                           # 统计缺失值
df.dropna()                                 # 删除缺失值
df['列'].fillna(值)                         # 填充缺失值
df['列'].fillna(df['列'].mean())           # 用均值填充

# 重复值
df.duplicated().sum()                       # 统计重复值
df.drop_duplicates()                        # 删除重复值

# 异常值（IQR方法）
Q1 = df['列'].quantile(0.25)
Q3 = df['列'].quantile(0.75)
IQR = Q3 - Q1
outliers = df[(df['列'] < Q1-1.5*IQR) | (df['列'] > Q3+1.5*IQR)]
\end{lstlisting}

\subsection{数据转换}

\begin{lstlisting}[caption={转换数据}]
# 类型转换
df['列'].astype(int)
df['列'].astype(str)
pd.to_datetime(df['日期列'])

# 字符串处理
df['列'].str.strip()                        # 去空格
df['列'].str.replace('旧', '新')            # 替换
df['列'].str.contains('关键词')             # 包含判断

# 应用函数
df['新列'] = df['列'].apply(函数)
\end{lstlisting}

\subsection{分组统计}

\begin{lstlisting}[caption={分组统计}]
# 简单分组
df.groupby('分组列')['统计列'].mean()

# 多列多统计
df.groupby(['列1', '列2']).agg({
    '列3': ['mean', 'sum'],
    '列4': 'count'
})
\end{lstlisting}

\end{document}
